\documentclass[11pt,a4paper]{article}

\usepackage[margin=1.8cm]{geometry}
\usepackage[T1]{fontenc}
\usepackage[utf8]{inputenc}
\usepackage{array}
\usepackage{enumitem}
\usepackage{hyperref}
\usepackage{xcolor}

\hypersetup{
  colorlinks=true,
  urlcolor=blue,
  linkcolor=black
}

\setlength{\parindent}{0pt}
\setlength{\parskip}{4pt}
\setlist[itemize]{leftmargin=1.2em,itemsep=2pt,topsep=2pt}

\title{\vspace{-1.5em}Group 3 Review of Group 5 DMP and Published Artefacts}
\author{Group 3}
\date{June 2026}

\begin{document}
\maketitle
\vspace{-1.5em}

\textbf{Review target:} Group 5, \url{https://test.researchdata.tuwien.ac.at/records/v888d-r0r44}

\textbf{Project:} Predicting Urban Traffic Congestion Levels using Multi-Layer Perceptrons on SDCC SCOOT Data

\section*{Summary Ratings}

\begin{center}
\begin{tabular}{|p{0.58\linewidth}|c|p{0.22\linewidth}|}
\hline
\textbf{Review dimension} & \textbf{Rating} & \textbf{Short note} \\
\hline
(a) Completeness of all FWF DMP fields & 3/3 & All main DMP sections are present. \\
\hline
(b) Quality and specificity of data descriptions & 2/3 & Data present and accessible; no semantic annotations on DBRepo columns. \\
\hline
(c) Correctness of metadata standard artefacts & 2/3 & Artefacts present; one README placeholder remains. \\
\hline
(d) maDMP technical validity & 2/3 & JSON is valid; stricter schema check recommended. \\
\hline
(e) Consistency between DMP text and published artefacts & 2/3 & Main links resolve; minor documentation mismatch. \\
\hline
\end{tabular}
\end{center}

\section*{(a) Completeness of all FWF DMP Fields -- 3/3}

The DMP record is published and openly accessible in TUWRD, and it contains both the human-readable DMP PDF and the machine-actionable maDMP JSON. The DMP follows the FWF DMP structure and includes administrative information, data responsibilities, reused and produced datasets, metadata and documentation, data quality, storage and backup, publication and preservation, legal aspects, and ethical aspects.

The DMP also explains the SDCC SCOOT traffic dataset, DBRepo reuse, GitHub and Zenodo software publication, TUWRD model and generated-output deposits, and the planned metadata standards. Overall, the DMP fields are complete enough for review, so this dimension receives the highest rating.

\section*{(b) Quality and Specificity of Data Descriptions -- 2/3}

The DBRepo database (SDCC Traffic Flow Experiment 2022, \url{https://test.dbrepo.tuwien.ac.at/database/4c334c86-77a1-44b7-a649-0588a5362a06}) is publicly accessible and contains three tables: \texttt{TrafficMeasurements} (1,048,575 rows of 15-minute traffic flow and congestion metrics), \texttt{TrafficSites} (physical sensor locations), and \texttt{Calendar} (date-to-day mapping). Three descriptively named SQL views are also present: \texttt{v\_peak\_hour\_measurements}, \texttt{v\_weekday\_measurements}, and \texttt{v\_measurements\_enriched}, which expose useful ML-relevant subsets of the data.

However, the \texttt{TrafficMeasurements} schema shows empty \textit{Concept} and \textit{Measurement Unit} fields for all columns -- no ontology concept URIs or unit of measurement annotations (such as QUDT or SI Digital Framework references) have been assigned to any column. The DMP text also does not describe the DBRepo schema or column-level metadata in detail. Because the data is present and findable but lacks semantic column-level annotations that would make it fully FAIR-compliant, this dimension is rated 2/3.

\section*{(c) Correctness of Metadata Standard Artefacts -- 2/3}

All five expected metadata artefacts (RO-Crate, CodeMeta, Croissant, FAIR4ML, and Model Card) are present in the repository and provide high-quality documentation. The RO-Crate (\texttt{ro-crate-metadata.json}) correctly maps the relationships between datasets, software, and evaluation outputs. The CodeMeta file includes detailed dependency tracking (e.g., TensorFlow, DBRepo) and versioning. The Croissant record defines the schema for the 3rd Normal Form tables and, notably, includes QUDT unit URIs (e.g., percent, count) for numeric attributes. The Model Card and FAIR4ML metadata provide comprehensive evaluation metrics and usage guidelines.

The main issue found during the check is that the README still contains a placeholder for the TUWRD model DOI, even though the model deposit is available through the DMP/TUWRD record. Additionally, while the metadata files are technically correct, the semantic unit mappings provided in the Croissant file are not yet reflected in the DBRepo database UI (see Section b). Because of these consistency issues, this dimension is rated 2/3.

\section*{(d) maDMP Technical Validity -- 2/3}

The maDMP file is syntactically valid JSON and contains an RDA-style top-level \texttt{dmp} object. It includes contact information, contributors, cost, created and modified dates, project information, ethical issue status, and five dataset entries with distributions.

The maDMP is therefore technically usable at the JSON level. However, some dataset entries are not very easy to inspect with simple JSON access because clear human-readable titles are not exposed consistently. A strict RDA schema validation should still be run before the final submission, so this dimension is rated 2/3.

\section*{(e) Consistency Between DMP Text and Published Artefacts -- 2/3}

The main published artefacts mentioned in the DMP could be found. The GitHub repository, Zenodo records, and TUWRD deposits for models and outputs are all consistent with the DMP's claims. The DBRepo database is publicly accessible and contains the expected tables and views.

However, a significant consistency gap exists between the metadata documentation and the actual repository state: the Croissant metadata file correctly identifies semantic units using QUDT, but these are missing from the DBRepo database UI. This discrepancy, combined with the README placeholder issue, leads to a rating of 2/3 for consistency.

\end{document}